% !TEX TS-program = pdflatex
% !TEX encoding = UTF-8 Unicode

% This is a simple template for a LaTeX document using the "article" class.
% See "book", "report", "letter" for other types of document.

\documentclass[11pt]{article} % use larger type; default would be 10pt
\usepackage[english,russian]{babel}
\usepackage{amsmath}
\usepackage[utf8]{inputenc} % set input encoding (not needed with XeLaTeX)

%%% Examples of Article customizations
% These packages are optional, depending whether you want the features they provide.
% See the LaTeX Companion or other references for full information.

%%% PAGE DIMENSIONS
\usepackage{geometry} % to change the page dimensions
\geometry{a4paper} % or letterpaper (US) or a5paper or....
\geometry{margin=0.5in} % for example, change the margins to 2 inches all round
% \geometry{landscape} % set up the page for landscape
%   read geometry.pdf for detailed page layout information

\usepackage{graphicx} % support the \includegraphics command and options

% \usepackage[parfill]{parskip} % Activate to begin paragraphs with an empty line rather than an indent

%%% PACKAGES
\usepackage{booktabs} % for much better looking tables
\usepackage{array} % for better arrays (eg matrices) in maths
\usepackage{paralist} % very flexible & customisable lists (eg. enumerate/itemize, etc.)
\usepackage{verbatim} % adds environment for commenting out blocks of text & for better verbatim
\usepackage{subfig} % make it possible to include more than one captioned figure/table in a single float
% These packages are all incorporated in the memoir class to one degree or another...

%%% HEADERS & FOOTERS
\usepackage{fancyhdr} % This should be set AFTER setting up the page geometry
\pagestyle{fancy} % options: empty , plain , fancy
\renewcommand{\headrulewidth}{0pt} % customise the layout...
\lhead{}\chead{}\rhead{}
\lfoot{}\cfoot{\thepage}\rfoot{}

%%% SECTION TITLE APPEARANCE
\usepackage{sectsty}
\allsectionsfont{\sffamily\mdseries\upshape} % (See the fntguide.pdf for font help)
% (This matches ConTeXt defaults)

%%% ToC (table of contents) APPEARANCE
\usepackage[nottoc,notlof,notlot]{tocbibind} % Put the bibliography in the ToC
\usepackage[titles,subfigure]{tocloft} % Alter the style of the Table of Contents
\renewcommand{\cftsecfont}{\rmfamily\mdseries\upshape}
\renewcommand{\cftsecpagefont}{\rmfamily\mdseries\upshape} % No bold!

%%% END Article customizations

%%% The "real" document content comes below...

\usepackage{titlesec}

% Настройка формата заголовка секции: жирный, обычный размер
\titleformat{\section}
  {\normalfont\Large\bfseries} % Шрифт: размер, жирный
  {\thesection}                % Нумерация
  {1em}                       % Расстояние
  {}           

\title{\textbf{Пересечение двух отрезков в пространстве}}
\author{Бирюков Вячеслав Геннадиевич}
%\date{} % Activate to display a given date or no date (if empty),
         % otherwise the current date is printed 

\begin{document}
\maketitle

\section{Постановка задачи}

Есть класс трёхмерного вектора:
\begin{verbatim}
class Vector3D
{
double X;
double Y;
double Z;
...

}
\end{verbatim}
и трёхмерного отрезка, заданного двумя Vector3D:

\begin{verbatim}
class Segment3D
{
Vector3D start;
Vector3D end;
...

}
\end{verbatim}
Требуется написать функцию Intersect, которая будет находить точку Vector3D пересечения двух
заданных на вход Segment3D. В классы Vector3D и Segment3D можно добавлять любые методы.


\section{Метод решения задачи}

Для решения поставенной в \textbf{п.1} задачи воспользуемся параметрическим уравнением отрезка прямой, соединяющего точки (концы векторов) $\mathbf{v}_{1}(x_{1}, y_{1}, z_{1})$ и $\mathbf{v}_{2}(x_{2}, y_{2}, z_{2})$:
\begin{equation}
	\mathbf{p}\left(t\right) = \mathbf{v}_{1} + t\left(\mathbf{v}_{2} - \mathbf{v}_{1}\right), \label{eq:1}
\end{equation}
где $t \in \left[0; 1\right]$.

Пусть параметрически с помощью формулы (\ref{eq:1}) заданы два отрезка:
\begin{eqnarray}
	\mathbf{p}\left(t\right) = \mathbf{v}_{1} + t\left(\mathbf{v}_{2} - \mathbf{v}_{1}\right), \label{eq:2} \\
	\mathbf{q}\left(t\right) = \mathbf{u}_{1} + t\left(\mathbf{u}_{2} - \mathbf{u}_{1}\right) \label{eq:3}. 
\end{eqnarray}
Для того чтобы эти отрезки пересекались, т.е. имели общую точку нужно потребовать для некоторых значений $t_{1}$ и $t_{2}$ выполнение следующего равенства:
\begin{equation}
	\mathbf{p}\left(t_{1}\right) = \mathbf{q}\left(t_{2}\right). \label{eq:4}
\end{equation}
С учетом соотношений (\ref{eq:2}) и (\ref{eq:3}) равенство (\ref{eq:4}) запишется в виде
\begin{equation}
	 \mathbf{v}_{1} + t\left(\mathbf{v}_{2} - \mathbf{v}_{1}\right) = 
	 \mathbf{u}_{1} + t\left(\mathbf{u}_{2} - \mathbf{u}_{1}\right). \label{eq:5}
\end{equation} 
Введем обозначения:
\begin{eqnarray} 
	\Delta \mathbf{v} = \mathbf{v_{2}} - \mathbf{v_{1}}, \label{eq:6}\\
	\Delta \mathbf{u} = \mathbf{u_{2}} - \mathbf{u_{1}}, \label{eq:7} \\
        \mathbf{w} = \mathbf{u_{1}} - \mathbf{v_{1}}. \label{eq:8}
\end{eqnarray}
Тогда уравнение (\ref{eq:5}) можно представить в виде:
\begin{equation}
	Ax=B, \label{eq:9}
\end{equation}
где
$$
	A = 
	\begin{pmatrix}
		\Delta v_{x} & -\Delta u_{x} \\
		\Delta v_{y} & -\Delta u_{y} \\
		\Delta v_{z} & -\Delta u_{z} \\				
	\end{pmatrix},
	x = 
	\begin{pmatrix}
	t_{1} \\ t_{2}
	\end{pmatrix},
	A = 
	\begin{pmatrix}
		w_{x}  \\
		w_{y}  \\
		w_{z}  \\				
	\end{pmatrix}.		
$$

Система линейных алгебраических уравнений (\ref{eq:9}) является переопределенной, т.е. количество уравнений больше чем количество неизветстных. К такого рода системам, в частности, приводит применение метода наименьших квадратов. Для решения таких систем разработаны специальные методы, напрмер: non-negative least squares (NNLS), псевдообратные 
матрицы Мура-Пенроуза, методы с использованием SVD-разложения и д.р. В некоторых случаях для обеспечения 
устойчивости решения применяют методы решения некорректно поставленных задач, в частности, метод регуляризации Тихонова. \footnote{Мной были реализрваны эти и многие другие методы в рамках работы над проектами ЯМК (прибор ядерно-магнитного каротажа) и 8АД (прибор акустического кроссдипольного каротажа).}

Для простоты вычислений воспользуемся наиболее простым подходом, также используемым при решении задач методом наименьших квадратов. Домножим обе части уравнения (\ref{eq:9}) на $A^{T}$ - транспонированную матрицу. Получим совместную систему линейных алгебраических уравнений 2-го порядка:
\begin{equation}
	A^{T}Ax=A^{T}B, \label{eq:10}
\end{equation}  
где 
$$
	A^{T}A = 
	\begin{pmatrix}
		|\Delta \mathbf{v}|^{2} & -\Delta \mathbf{v} \cdot \Delta \mathbf{u}  \\
		-\Delta \mathbf{v} \cdot \Delta \mathbf{u} & |\Delta \mathbf{u}|^{2}			
	\end{pmatrix},
	A^{T}B = 
	\begin{pmatrix}
		\Delta \mathbf{v} \cdot \Delta \mathbf{w}  \\
		-\Delta \mathbf{u} \cdot \Delta \mathbf{w} 
	\end{pmatrix},
$$
где <<$\cdot$>> обозначает скалярное произведение, $|\cdot|^{2}$ - квадрат модуля вектора.

Решение системы уравнений (\ref{eq:10}) легко найти, напрмер, методом Крамера:
\begin{equation}
	t_{1} = \frac{D_{1}}{D}, 	t_{2} = \frac{D_{2}}{D},  \label{eq:11}
\end{equation}
где
\begin{gather*}	
	D = |\Delta \mathbf{v}|^{2}|\Delta \mathbf{u}|^{2} - \left(\Delta \mathbf{v} \cdot \Delta \mathbf{u}\right)^{2}, \\		
	D_{1} = |\Delta \mathbf{u}|^{2}\left(\Delta \mathbf{v} \cdot \Delta \mathbf{w}\right) 
	- \left(\Delta \mathbf{v} \cdot \Delta \mathbf{u}\right)\left(\Delta \mathbf{u} \cdot \Delta \mathbf{w}\right), \\
	D_{2} = |\Delta \mathbf{v}|^{2}\left(\Delta \mathbf{u} \cdot \Delta \mathbf{w}\right) 
	- \left(\Delta \mathbf{v} \cdot \Delta \mathbf{u}\right)\left(\Delta \mathbf{v} \cdot \Delta \mathbf{w}\right).	
\end{gather*}

Найти координаты точки пересечения можно путем подстановки (\ref{eq:11}) в (\ref{eq:2}) или (\ref{eq:3}).

Сделаем несколько замечаний:

\textbf{Замечание 1.} Опредилитель $D = |\Delta \mathbf{v}|^{2}|\Delta \mathbf{u}|^{2} - \left(\Delta \mathbf{v} \cdot \Delta \mathbf{u}\right)^{2}$ будет равен нулю для параллельных векторов $\Delta \mathbf{v} || \Delta \mathbf{u}$. Это следует из определения скалярного произведения: 
	$\mathbf{a} \cdot \mathbf{b} = |\mathbf{a}||\mathbf{b}|\cos(\mathbf{a}, \mathbf{b})$. Следовательно, при реализации алгоритма нужно организовать проверку на близость к нулю $D$ и обработать эту ситуацию отдельно.  

\textbf{Замечание 2.} Условия $t_{1} \in \left[0; 1\right]$ и $t_{2} \in \left[0; 1\right]$ могут не выполниться: точка пересечения лежит не внутри заданных отрезков. Также эту ситуацию следует обаботать отдельно.

\textbf{Замечание 3.} При программной реализации алгоритма следует организовать проверку выполнения равества (\ref{eq:4}):
$$
	|\mathbf{p}\left(t_{1}\right) - \mathbf{q}\left(t_{2}\right)| < \epsilon,
$$
если такое неравенство не выполняет отрезки не пересекаются.

\section{Реализация классов Vector3D и Segment3D}

\subsection{Vector3D.h}
\begin{verbatim}
#include<cmath>

#pragma once
class Vector3D
{
public:

	double X;
	double Y;
	double Z;

	Vector3D(double x = 0, double y = 0, double z = 0)
		: X(x), Y(y), Z(z)
	{ }

	Vector3D operator+(const Vector3D& other) const;
	Vector3D operator-(const Vector3D& other) const;
	Vector3D operator*(double scalar) const;
	double dot(const Vector3D& other) const;
	Vector3D cross(const Vector3D& other) const;
	double length() const;
	Vector3D normalized() const;
};
\end{verbatim}

\subsection{Vector3D.cpp}
\begin{verbatim}
#include "Vector3D.h"

Vector3D Vector3D::operator+(const Vector3D& other) const 
{
    return Vector3D(X + other.X, Y + other.Y, Z + other.Z);
}

Vector3D Vector3D::operator-(const Vector3D& other) const
{
	return Vector3D(X - other.X, Y - other.Y, Z - other.Z);
}

Vector3D Vector3D::operator*(double scalar) const
{
	return Vector3D(X * scalar, Y * scalar, Z * scalar);
}

double Vector3D::dot(const Vector3D& other) const
{
    return X * other.X + Y * other.Y + Z * other.Z;
}

Vector3D Vector3D::cross(const Vector3D& other) const 
{
	return Vector3D(
		Y * other.Z - Z * other.Y,
		Z * other.X - X * other.Z,
		X * other.Y - Y * other.X
	);
}

double Vector3D::length() const {
	return std::sqrt(X * X + Y * Y + Z * Z);
}

Vector3D Vector3D::normalized() const {
	double len = length();
	if (len < 1e-10) return Vector3D(0, 0, 0);
	return Vector3D(X / len, Y / len, Z / len);
}
\end{verbatim}


\subsection{Segment3D.h}
\begin{verbatim}
#include "Vector3D.h"


#pragma once
class Segment3D
{
public:
    Vector3D start;
    Vector3D end;

    Segment3D(const Vector3D& s, const Vector3D& e) : start(s), end(e) {}

    Vector3D direction() const;
    double length() const;     
};
\end{verbatim}

\subsection{Segment3D.cpp}
\begin{verbatim}
#include "Segment3D.h"

Vector3D Segment3D::direction() const 
{
	return end - start;
}

double Segment3D::length() const 
{
    return direction().length();
}
\end{verbatim}

\end{document}